% ID: SEL-D0004-OB-E02
% ESTADO: validada. S005, 2026-11-07; tres pasadas del mismo autor y revisión matemática, didáctica y visual.
\textbf{a) Continuidad y derivabilidad de f.}\par\nopagebreak[4]
El dominio es $\mathbb R$. Cada rama es continua y derivable en su intervalo abierto. Sólo hay que examinar $x=0$:
\[
\lim_{x\to0^-}(x^2+4x+1)=1,\quad
\lim_{x\to0^+}e^{-x}=1,\quad f(0)=1.
\]
Coinciden límites y valor; $f$ es \textbf{continua en todo $\mathbb R$}.

Para derivabilidad, comparamos derivadas laterales:
\[
f'_-(0)=\lim_{x\to0^-}(2x+4)=4,\qquad
f'_+(0)=\lim_{x\to0^+}(-e^{-x})=-1.
\]
Como $4\ne-1$, \textbf{no es derivable en 0}. En los demás puntos sí lo es:
\[
f'(x)=\begin{cases}2x+4&x<0,\\-e^{-x}&x>0.\end{cases}
\]
Comprobación directa con cocientes incrementales: a la izquierda, $[f(h)-1]/h=h+4\to4$; a la derecha, $(e^{-h}-1)/h\to-1$. La continuidad no garantiza derivabilidad.

\textbf{b) Parámetros de g.}\par\nopagebreak[4]
Un extremo relativo en $x=-1$ exige $g'(-1)=0$ y el paso por $(-1,3)$ exige $g(-1)=3$:
\[
g'(x)=3x^2+2bx,\qquad 3-2b=0\Rightarrow b=\frac32.
\]
\[
-1+b+c=3\Rightarrow c=4-b=\frac52.
\]
Por tanto, \boxed{b=3/2,\quad c=5/2}. Hay que comprobar que el punto estacionario sí es extremo: $g''(x)=6x+2b$ y $g''(-1)=-3<0$, luego es un máximo relativo estricto. Sustituyendo, $g(-1)=-1+3/2+5/2=3$ y $g'(-1)=3-3=0$.
